% LaTeX master Datei(en) zusammengestellt von Falk-Jonatan Strube zur Nutzung an der Hochschule für Technik und Wirtschaft Dresden: https://github.com/genericFJS/htw
\documentclass{scrreprt}
\gdef\pathtomaster{../_LaTeX_master}
%% HTW Dresden Corporate Design
%% Copyright (C) Falk-Jonatan Strube <falkjonatanstrube@gmail.com>, 2017
%%% === Anfang Paket htwscr ===
\makeatletter
\ifcsname KOMAoption\endcsname%
\KOMAoption{footheight}{11mm}	% mehr Platz für Wasserzeichen o.ä. in der Fußzeile
\recalctypearea	% s.o.
\RequirePackage{geometry}
\geometry{top=25mm, bottom=25mm, headsep=10mm, inner=25mm, outer=25mm, footskip=15mm, marginparsep=2mm}
\RequirePackage{helvet}	% Corporate Design schreibt Arial/Helvetica vor
\renewcommand{\familydefault}{\sfdefault}	% s.o.
\RequirePackage{graphicx}	% für Grafiken
\RequirePackage{xcolor}		% für Farben
\definecolor{htworange}{RGB}{249,155,28}	% Corporate Design orange CMYK(0/45/100/0)
\definecolor{htwblue}{RGB}{0,116,188}			% Corporate Design blau CMYK(90/50/0/0)
\definecolor{htwgrey}{RGB}{128,128,128}		% Corporate Design gray 60% Schwarz [Annahme: CMYK(0,0,0,60)]
\RequirePackage{ifthen,xifthen}	% für Bedingungen/Optionen
\RequirePackage{etoolbox}				% s.o.

%%% Format des HTW-Headers (wird hier als Texteelement verwenden) nach dem Style der PowerPoint-Vorlagen. Weiterhin inspiriert durch das Paket tudscr (Corporate Design der TU Dresden)
%%% Dateipfad des Logos abhängig von der Verwendung als Paket/Input
\ifdefined\@isHTWCDpackage%
	\gdef\@pathtoimage{HTW}%
\else%
	\ifdefined\pathtomaster%
		\gdef\@pathtoimage{\pathtomaster /HTW}%
	\else 
		\gdef\@pathtoimage{HTW}%
	\fi%
\fi%
\newcommand{\pathtologo}[1]{\gdef\@pathtoimage{#1}}
%%% Berechnung der Größe des Logos
\newlength{\htwLogoHeight}
\setlength{\htwLogoHeight}{12mm}	% Mindesthöhe des Logos laut Corporate Design: 11mm
\newlength{\@htwLogoWidth}
\AtBeginDocument{
\def\printLogo{\includegraphics[height=\htwLogoHeight]{\@pathtoimage}}
\setlength\@htwLogoWidth{\widthof{\printLogo}}
}
%%% Definition des Header
\newcommand{\@htwheader}{{\noindent%
%%% Header ist abhängig von dem Vorhandensein von department, institute und chair
\newcommand{\hasDepInsCh}{\gdef\@hasDepInsCh{}}%
\ifdefined\@department\hasDepInsCh\fi%
\ifdefined\@institute\hasDepInsCh\fi%
\ifdefined\@chair\hasDepInsCh\fi%
\begin{minipage}[b][\htwLogoHeight]{\dimexpr \textwidth - \@htwLogoWidth - 2mm \relax}% 2mm Abstand zwischen Logo und Linie
{\small{\bfseries%
\ifdefined\@faculty\@faculty{}\fi} %
\strut}%
{\color{htworange}\hrule height 1pt}%
\vspace{.1mm}%
\end{minipage}
\hfill%
\begin{minipage}[b][\htwLogoHeight]{\@htwLogoWidth}
\includegraphics[height=\htwLogoHeight]{\@pathtoimage}
\end{minipage}
\ifdefined\@hasDepInsCh%
\\%
\begin{minipage}{\textwidth}
\small%
\ifdefined\@department\@department{}\fi%
\ifdefined\@institute\ifdefined\@department,~\fi \@institute\fi%
\ifdefined\@chair\ifdefined\@department,~\else\ifdefined\@institute,~\fi\fi\@chair\fi\strut
%{\color{htworange}\hrule height 0.5pt}
\end{minipage}
\fi
}}

%%% Definition der Fußzeilen
\usepackage{scrlayer-scrpage}
\@ifclassloaded{scrreprt}{%
\renewcommand*{\partpagestyle}{scrheadings}
\renewcommand*{\chapterpagestyle}{scrheadings}%
}{}
\pagestyle{scrheadings}
\cfoot{\vfill\pagemark}
\chead{}

%%% Einstellungen
\newcommand{\simpledate}{\gdef\@simpledate{}}		% Datum nicht mit "Eingereicht am:" Vorsatz
\newcommand{\nouppercase}{\gdef\@nouppercase{}}	% Titel/Parts/Chapters usw. nicht in Großbuchstaben

%%% Einiges erst zu Beginn des Dokumentes ausführen, um Einstellungen zu berücksichtigen
\AtBeginDocument{
\renewcommand\SS{SS}	% fix big ß
%%% Sektionen usw. ggf. in Großbuchstaben:
\providecommand{\chapterlinesformat}[3]{}	% compatibility
\providecommand{\chapterlineswithprefixformat}[3]{}	% compatibility
\providecommand{\sectionlinesformat}[4]{}	% compatibility
\ifdefined\@nouppercase\relax\else
\renewcommand\sectionlinesformat[4]{\@hangfrom{\hskip#2 #3}{\MakeUppercase{#4}}}
\@ifclassloaded{scrreprt}{%
\renewcommand\chapterlinesformat[3]{\@hangfrom{#2}{\MakeUppercase{#3}}}
\renewcommand\chapterlineswithprefixformat[3]{\MakeUppercase{#2#3}}%
}{}
\fi
%%% Anpassung der Part-Seite
\renewcommand*{\partheadstartvskip}{\@htwheader\vspace*{30mm}}
\renewcommand*{\raggedpart}{\raggedright}
\renewcommand*{\partformat}{\huge\ifdefined\@nouppercase \partname~\thepart\else\MakeUppercase{\partname~\thepart}\fi.}
\ifdefined\@nouppercase\setkomafont{part}{\Huge}\else\setkomafont{part}{\Huge\MakeUppercase}\fi
}

\newcommand{\bilingual}[3]{\ifdefined #1 \let #1\undefined\fi\newcommand{#1}{\iflanguage{english}{#3}{#2}}}	% Übersetzung von einigen Variablen in mehrere Sprachen

\newcommand{\checkThesis}[2]{\ifthenelse{\equal{#1}{diss}\OR\equal{#1}{doctoral}\OR\equal{#1}{phd}}{
\gdef#2{\dissertationname}}{
\ifthenelse{\equal{#1}{diploma}}{
\gdef#2{\diplomathesisname}}{
\ifthenelse{\equal{#1}{master}}{
\gdef#2{\masterthesisname}}{
\ifthenelse{\equal{#1}{bachelor}}{
\gdef#2{\bachelorthesisname}}{
\ifthenelse{\equal{#1}{student}}{
\gdef#2{\studentthesisname}}{
\ifthenelse{\equal{#1}{evidence}}{
\gdef#2{\studentresearchname}}{
\ifthenelse{\equal{#1}{project}}{
\gdef#2{\projectpapername}}{
\ifthenelse{\equal{#1}{seminar}}{
\gdef#2{\seminarpapername}}{
\ifthenelse{\equal{#1}{term}}{
\gdef#2{\termpapername}}{
\ifthenelse{\equal{#1}{research}}{
\gdef#2{\researchname}}{
\ifthenelse{\equal{#1}{log}}{
\gdef#2{\logname}}{
\ifthenelse{\equal{#1}{report}}{
\gdef#2{\reportname}}{
\ifthenelse{\equal{#1}{internship}}{
\gdef#2{\internshipname}}{
\ifthenelse{\equal{#1}{lecture}\OR\equal{#1}{lesson}\OR\equal{#1}{pract}}{
\bilingual{\professorname}{Vorlesung von}{Lecture by}
\simpledate
\gdef\@islecture{\relax}
\ifthenelse{\equal{#1}{lecture}}{\gdef#2{\lecturename}}{
\ifthenelse{\equal{#1}{lesson}}{\gdef#2{\lessonname}}{
\gdef#2{\practname}
}}}{
\gdef#2{#1}}}}}}}}}}}}}}}}

%%% Textfragmente der Titelseite in deutsch und englisch
\bilingual{\dateofbirthname}{Geboren am:}{Born on:}
\bilingual{\placeofbirthname}{in}{in}
\bilingual{\coursename}{Studiengang:}{Course:}
\bilingual{\disciplinename}{Studienrichtung:}{Discipline:}
\bilingual{\matriculationnumbername}{Matrikelnummer:}{Matriculation number:}
\bilingual{\matriculationyearname}{Immatrikulationsjahr}{Matriculation year:}
\bilingual{\graduationname}{zur Erlangung des akademischen Grades}{to achieve the academic degree}
\bilingual{\refereename}{Gutachter}{Referee}
\bilingual{\advisorname}{Fachreferent}{Advisor}
\bilingual{\supervisorname}{Betreuer}{Supervisor}
\bilingual{\professorname}{Betreuender Hochschullehrer}{Supervising professor}
\bilingual{\datename}{Eingereicht am:}{Submitted on:}
%% Thesis Typen (inspiriert von dem Paket tudscr):
\bilingual{\dissertationname}{Dissertation}{Dissertation} % diss, doctoral, phd 
\bilingual{\diplomathesisname}{Diplomarbeit}{Diploma Thesis} % diploma 
\bilingual{\masterthesisname}{Master-Arbeit}{Master Thesis} % master 
\bilingual{\bachelorthesisname}{Bachelor-Arbeit}{Bachelor Thesis} % bachelor 
\bilingual{\studentthesisname}{Studienarbeit}{Student Thesis} % student 
\bilingual{\studentresearchname}{Großer Beleg}{Student Research Project} % evidence 
\bilingual{\projectpapername}{Projektarbeit}{Project Paper} % project
\bilingual{\seminarpapername}{Seminararbeit}{Seminar Paper} % seminar
\bilingual{\termpapername}{Hausarbeit}{Term Paper} % term
\bilingual{\researchname}{Forschungsbericht}{Research Report} % research 
\bilingual{\logname}{Protokoll}{Log} % log 
\bilingual{\reportname}{Bericht}{Report} % report
\bilingual{\internshipname}{Praktikumsbericht}{Internship Report} % internship 
\bilingual{\lecturename}{Vorlesungsmitschrift}{Lecture Notes} % lecture 
\bilingual{\lectureauthorname}{Mitschrift von}{Notes by} % lecture  
\bilingual{\lessonname}{\"Ubungsmitschrift}{Lesson Notes} % lecture  
\bilingual{\practname}{Praktikumsmitschrift}{Pracital Course Notes} % lecture 
%%% Variablen für die Titelseite (inspiriert von dem Paket tudscr):
\renewcommand{\subject}[1]{\checkThesis{#1}{\@subject}}
\newcommand{\faculty}[1]{\gdef\@faculty{#1}}
\newcommand{\department}[1]{\gdef\@department{#1}}
\newcommand{\institute}[1]{\gdef\@institute{#1}}
\newcommand{\chair}[1]{\gdef\@chair{#1}}
\newcommand{\authormore}[1]{\gdef\@authormore{#1}}
\newcommand{\dateofbirth}[1]{\gdef\@dateofbirth{#1}}
\newcommand{\placeofbirth}[1]{\gdef\@placeofbirth{#1}}
\newcommand{\course}[1]{\gdef\@course{#1}}
\newcommand{\discipline}[1]{\gdef\@discipline{#1}}
\newcommand{\matriculationnumber}[1]{\gdef\@matriculationnumber{#1}}
\newcommand{\matriculationyear}[1]{\gdef\@matriculationyear{#1}}
\newcommand{\thesis}[1]{\checkThesis{#1}{\@thesis}}
\newcommand{\graduation}[2][]{\gdef\@graduation{#2}\gdef\@graduationshort{#1}}
\newcommand{\supervisor}[1]{\gdef\@supervisor{#1}}
\newcommand{\referee}[1]{\gdef\@referee{#1}}
\newcommand{\advisor}[1]{\gdef\@advisor{#1}}
\newcommand{\professor}[1]{\gdef\@professor{#1}}
%%% Titelseite (inspiriert von dem Paket tudscr):
\renewcommand*{\maketitle}{%
\begin{titlepage}%
\begin{flushleft}
\renewcommand{\thefootnote}{$\star$}%
%%% Kopfzeile HTW-Style:
\@htwheader\\[13mm]
%%% Die Zeilen für titlehead, subject, title, subtitle und author sind immer vorhanden (spacing)
{\noindent\@titlehead{}\strut}\\[3mm]%
{\Large\@subject{}\strut}\\[8mm]%
{\Huge\bfseries%
\ifdefined\@nouppercase\@title{}\else\MakeUppercase{\@title{}}\fi%
\strut}\\
{\Large\bfseries%
\ifdefined\@nouppercase\@subtitle{}\else\MakeUppercase{\@subtitle{}}\fi%
\strut}\\[11mm]
\ifdefined\@islecture\strut\else{\Large\@author{}\strut}\fi % Autor wird bei Vorlesungsmitschrift unten angezeigt
%%% Weitere Felder werden nur angezeigt, wenn sie definiert wurden
\ifdefined\@authormore\ifdefined\@islecture\else\\\fi{\@authormore{}\strut}\fi
\ifdefined\@dateofbirth\\\dateofbirthname{} \@dateofbirth{} \ifdefined\@placeofbirth\placeofbirthname{} \@placeofbirth{}\strut\fi\fi
\ifdefined\@course\\\coursename{} \@course{}\strut\fi
\ifdefined\@discipline\\\disciplinename{} \@discipline{}\strut\fi
\ifdefined\@matriculationnumber\\\matriculationnumbername{} \@matriculationnumber{}\strut\fi
\ifdefined\@matriculationyear\\\matriculationyearname{} \@matriculationyear{}\strut\fi
\ifdefined\@thesis\\[10mm]{\LARGE\bfseries
\ifdefined\@nouppercase\@thesis{}\else\MakeUppercase{\@thesis{}}\fi
\strut}
\ifdefined\@graduation\\[2mm]{\graduationname{}\strut}
\\[5mm]{\LARGE\bfseries
\ifdefined\@nouppercase\relax\else\MakeUppercase{\fi%
\@graduation{} \ifdefined\@graduationshort(\@graduationshort{})\fi
\ifdefined\@nouppercase\relax\else}\fi\strut}\fi\fi
\vfill\noindent
\ifdefined\@islecture\\[1mm]{\lectureauthorname{}\strut}\\[.7mm]{\Large\@author{}\strut}\fi
\ifdefined\@referee{\refereename{}\strut}\\[.7mm]{\Large\@referee{}\strut}\fi
\ifdefined\@advisor\\[1mm]{\advisorname{}\strut}\\[.7mm]{\Large\@advisor{}\strut}\fi
\ifdefined\@supervisor\\[1mm]{\supervisorname{}\strut}\\[.7mm]{\Large\@supervisor{}\strut}\fi
\ifdefined\@professor\\[1mm]{\professorname{}\strut}\\[.7mm]{\Large\@professor{}\strut}\fi
\ifdefempty{\@date}{}{\ifdefined\@simpledate\\[6mm]\relax\else\\[3mm]\datename{}~\fi\@date{}\strut}
\ifdefempty{\@publishers}{}{\\[4.5mm]{\Large\@publishers{}\strut}}%
\renewcommand\footnoterule{}
\@thanks\global\let\@thanks\@empty
\end{flushleft}
\end{titlepage}
\setcounter{footnote}{0}
}
\else%
\PackageError{htwcd}{htwcd needs a KOMA class!}%
\fi%
\makeatother
%%% === Ende Paket htwscr ===
%-- Nummerierung bis Subsubsection für Report
\setcounter{secnumdepth}{3}	% zählt auch subsubsection
\setcounter{tocdepth}{3}	% Inhaltsverzeichnis bis in subsubsection

% LuaLaTeX Unterscheidung
\usepackage{ifluatex}

% Input inkl. Umlaute, Silbentrennung
\ifluatex
	\usepackage{fontspec}
\else
	\usepackage[utf8]{inputenc}
	\usepackage[T1]{fontenc}
\fi
\usepackage[english,ngerman]{babel}
\usepackage{csquotes}	% Anführungszeichen
\usepackage{marvosym}
\usepackage{eurosym}

% Style-Aufhübschung
\usepackage{soul, color}	% Kapitälchen, Unterstrichen, Durchgestrichen usw. im Text
%\usepackage{titleref}
\usepackage[breakwords, fit]{truncate}	% Abschneiden von Sätzen
\renewcommand{\TruncateMarker}{\,…}

% Mathe usw.
\usepackage{amssymb}
\usepackage{amsthm}
\usepackage[fleqn,intlimits]{amsmath}	% fleqn: align-Umgebung rechtsbündig; intlimits: Integralgrenzen immer ober-/unterhalb
%\usepackage{mathtools} % u.a. schönere underbraces
\usepackage{xcolor}
\usepackage{esint}	% Schönere Integrale, \oiint vorhanden
\everymath=\expandafter{\the\everymath\displaystyle}	% Mathe Inhalte werden weniger verkleinert
\usepackage{wasysym}	% mehr Symbole, bspw \lightning
%\renewcommand{\int}{\int\limits}
%\usepackage{xfrac}	% mehr fracs: sfrac{}{}
\let\oldemptyset\emptyset	% schöneres emptyset
\let\emptyset\varnothing
%\RequirePackage{mathabx}	% mehr Symbole
\mathchardef\mhyphen="2D	% Hyphen in Math
\usepackage{mathtools}		% zum Beispiel {pmatrix}[r] für Formatierung der Matrix

% tikz usw.
\usepackage{tikz}
\usepackage{pgfplots}
\pgfplotsset{compat=1.11}	% Umgeht Fehlermeldung
\usetikzlibrary{graphs}
%\usetikzlibrary{through}	% ???
\usetikzlibrary{arrows}
\usetikzlibrary{arrows.meta}	% Pfeile verändern / vergrößern: \draw[-{>[scale=1.5]}] (-3,5) -> (-3,3);
\usetikzlibrary{automata,positioning} % Zeilenumbruch im Node node[align=center] {Text\\nächste Zeile} automata für Graphen
\usetikzlibrary{matrix}
\usetikzlibrary{patterns}	% Schraffierte Füllung
\usetikzlibrary{shapes.geometric}	% Polygon usw.
\tikzstyle{reverseclip}=[insert path={	% Inverser Clip \clip
	(current page.north east) --
	(current page.south east) --
	(current page.south west) --
	(current page.north west) --
	(current page.north east)}
% Nutzen: 
%\begin{tikzpicture}[remember picture]
%\begin{scope}
%\begin{pgfinterruptboundingbox}
%\draw [clip] DIE FLÄCHE, IN DER OBJEKT NICHT ERSCHEINEN SOLL [reverseclip];
%\end{pgfinterruptboundingbox}
%\draw DAS OBJEKT;
%\end{scope}
%\end{tikzpicture}
]	% Achtung: dafür muss doppelt kompliert werden!
\usepackage{graphpap}	% Grid für Graphen
\tikzset{every state/.style={inner sep=2pt, minimum size=2em}}
\usetikzlibrary{mindmap, backgrounds}
%\usepackage{tikz-uml}	% braucht Dateien: http://perso.ensta-paristech.fr/~kielbasi/tikzuml/

% Tabular
\usepackage{longtable}	% Große Tabellen über mehrere Seiten
\usepackage{multirow}	% Multirow/-column: \multirow{2[Anzahl der Zeilen]}{*[Format]}{Test[Inhalt]} oder \multicolumn{7[Anzahl der Reihen]}{|c|[Format]}{Test2[Inhalt]}
\renewcommand{\arraystretch}{1.3} % Tabellenlinien nicht zu dicht
\usepackage{colortbl}
\arrayrulecolor{gray}	% heller Tabellenlinien
\usepackage{array}	% für folgende 3 Zeilen (für Spalten fester breite mit entsprechender Ausrichtung):
\newcolumntype{L}[1]{>{\raggedright\let\newline\\\arraybackslash\hspace{0pt}}m{\dimexpr#1\columnwidth-2\tabcolsep-1.5\arrayrulewidth}}
\newcolumntype{C}[1]{>{\centering\let\newline\\\arraybackslash\hspace{0pt}}m{\dimexpr#1\columnwidth-2\tabcolsep-1.5\arrayrulewidth}}
\newcolumntype{R}[1]{>{\raggedleft\let\newline\\\arraybackslash\hspace{0pt}}m{\dimexpr#1\columnwidth-2\tabcolsep-1.5\arrayrulewidth}}
\usepackage{booktabs}	% Kompatibilität für array und \hrule
\usepackage{caption}	% Um auch unbeschriftete Captions mit \caption* zu machen

% Nützliches
\usepackage{verbatim}	% u.a. zum auskommentieren via \begin{comment} \end{comment}
\usepackage{tabto}	% Tabs: /tab zum nächsten Tab oder /tabto{.5 \CurrentLineWidth} zur Stelle in der Linie
\NumTabs{6}	% Anzahl von Tabs pro Zeile zum springen
\usepackage{listings} % Source-Code mit Tabs
\usepackage{lstautogobble} 
\usepackage{enumitem}	% Anpassung der enumerates
%\setlist[enumerate,1]{label=(\arabic*)}	% global andere Enum-Items
\renewcommand{\labelitemiii}{$\scriptscriptstyle ^\blacklozenge$} % global andere 3. Item-Aufzählungszeichen
\usepackage{letltxmacro} % neue Definiton von Grundbefehlen
% Nutzen:
%\LetLtxMacro{\oldemph}{\emph}
%\renewcommand{\emph}[1]{\oldemph{#1}}
\RequirePackage{xpatch}	% ua. Konkatenieren von Strings/Variablen (etoolbox)
\usepackage{xstring}	% String Operationen
\usepackage{minibox}	% Minibox anstatt \fbox{} für Boxen mit (manuellen) Zeilenumbruch
\usepackage{framed}		% \begin{framed} statt minibox/fbox mit automatischem Zeilenumbruch
\usepackage{microtype}	% Verfeinerte Hyphenation
\usepackage{import}		
\usepackage{standalone}	% Zusammen mit import: Einfügen von separaten Latex-Dateien (die schon header haben)

% Einrichtung von lst
%\renewcommand*{\ttdefault}{qcr}
\lstset{
basicstyle=\ttfamily, 
%basicstyle=\footnotesize\sffamily\color{black},
%mathescape=true, 
%escapeinside=^^, 
autogobble, 
tabsize=2,
frame=single,
rulecolor=\color{lightgray},
numbers=left,
numbersep=5pt,
numberstyle=\tiny\color{gray},
commentstyle=\color{gray},
keywordstyle=\color{htwblue},
stringstyle=\color{htworange},
morecomment=[l][\color{magenta}]{\#}
showspaces=false,
showstringspaces=false,
breaklines=true,
literate=%
    {Ö}{{\"O}}1
    {Ä}{{\"A}}1
    {Ü}{{\"U}}1
    {ß}{{\ss}}1
    {ü}{{\"u}}1
    {ä}{{\"a}}1
    {ö}{{\"o}}1
    {„}{{\glqq}}1
    {“}{{\grqq}}1
    {~}{{\textasciitilde}}1
}
\usepackage{scrhack} % Fehler umgehen
\def\ContinueLineNumber{\lstset{firstnumber=last}} % vor lstlisting. Zum wechsel zum nicht-kontinuierlichen muss wieder \StartLineAt1 eingegeben werden
\def\StartLineAt#1{\lstset{firstnumber=#1}} % vor lstlisting \StartLineAt30 eingeben, um bei Zeile 30 zu starten
\let\numberLineAt\StartLineAt

% BibTeX
\usepackage[
backend=bibtex,
bibencoding=ascii,
style=authoryear, citestyle=authoryear-ibid,
%doi=false,
%isbn=false,
%url=false
]{biblatex}	% BibTeX
\usepackage{makeidx}
%\makeglossary
%\makeindex
\usepackage{url}	% Darstellung von URLs

% Grafiken
\usepackage{graphicx}
\usepackage{epstopdf}	% eps-Vektorgrafiken einfügen
\usepackage{transparent}	% transparent nutzen: {\transparent{0.4} ...}
%\epstopdfsetup{outdir=\customDir}
% Prüft, ob Grafik existiert (mit \ifvalidimage{}{}) [Quelle: https://tex.stackexchange.com/a/99176]:
\makeatletter
\newif\ifgraphicexist
\catcode`\*=11
\newcommand\ifvalidimage[1]{%
    \begingroup
    \global\graphicexisttrue
    \let\input@path\Ginput@path
    \filename@parse{#1}%
    \ifx\filename@ext\relax
    \@for\Gin@temp:=\Gin@extensions\do{%
        \ifx\Gin@ext\relax
        \Gin@getbase\Gin@temp
        \fi}%
    \else
    \Gin@getbase{\Gin@sepdefault\filename@ext}%
    \ifx\Gin@ext\relax
    \global\graphicexistfalse
    \def\Gin@base{\filename@area\filename@base}%
    \edef\Gin@ext{\Gin@sepdefault\filename@ext}%
    \fi
    \fi
    \ifx\Gin@ext\relax
    \global\graphicexistfalse
    \else 
    \@ifundefined{Gin@rule@\Gin@ext}%
    {\global\graphicexistfalse}%
    {}%
    \fi  
    \ifx\Gin@ext\relax 
    \gdef\imageextension{unknown}%
    \else
    \xdef\imageextension{\Gin@ext}%
    \fi 
    \endgroup 
    \ifgraphicexist
    \expandafter \@firstoftwo
    \else
    \expandafter \@secondoftwo
    \fi 
} 
\catcode`\*=12
\makeatother
\LetLtxMacro{\forceincludegraphics}{\includegraphics}	% neuer Befehl für includegraphics
\renewcommand{\includegraphics}[2][]{	% altes includegraphics neu definieren, damit es auch nicht vorhandene einfügt
\ifvalidimage{#2}{
\forceincludegraphics[#1]{#2}
}{
\message{Achtung: Grafik wurde nicht gefunden: '#2'}
\minibox[frame]{
\textbf{\StrSubstitute{#2}{_}{\_}}%  \ifthenelse{\isempty{#1}}{}{\\\textit{#1}}
}}}

% pdf-Setup
\usepackage{pdfpages}
\usepackage[bookmarks,%
bookmarksopen=false,% Klappt die Bookmarks in Acrobat aus
colorlinks=true,%
linkcolor=black,%
citecolor=black,%
urlcolor=htworange,%
]{hyperref}

% Nach Dokumentenbeginn ausführen:
\makeatletter
\AtBeginDocument{	
	% pdf Informationen
	\hypersetup{
		pdfauthor={\@author},
		pdftitle={\@title},
		pdfproducer={LaTeX Layout erstellt von Falk-Jonatan Strube}
	}
	% signierte Fußzeile
	%\providecommand{\ofootname}{Mitschrift von}
	\providecommand{\pathtomaster}{../}
	\ofoot{%
	%\ifdefined\ofoottext
	%\else
	%\begin{minipage}{40mm}
	%\begin{flushright}
	%\textcolor{gray}{
	%\ofootname\\
	%Falk-Jonatan Strube\vspace*{-1mm}}
	%\end{flushright}
	%\end{minipage}
	%\fi%
	%\hspace*{1mm}%
	\begin{minipage}{10 mm}
	\includegraphics[height=9mm]{\pathtomaster/FJS-logo_gray}
	\end{minipage}%
	}
}
\makeatother



%%% ==================================================
%%% ===================== MACROS =====================
%%% ==================================================
% Horizontale Linie:
\newcommand{\HRule}[1][\medskipamount]{\par
  \vspace*{\dimexpr-\parskip-\baselineskip+#1}
  \noindent\rule[0.2ex]{\linewidth}{0.2mm}\par
  \vspace*{\dimexpr-\parskip-.5\baselineskip+#1}}
% Gestrichelte horizontale Linie:
\RequirePackage{dashrule}
\newcommand{\HDRule}[1][\medskipamount]{\par
  \vspace*{\dimexpr-\parskip-\baselineskip+#1}
  \noindent\hdashrule[0.2ex]{\linewidth}{0.2mm}{1mm} \par
  \vspace*{\dimexpr-\parskip-.5\baselineskip+#1}}
% Mathe in Anführungszeichen:
\newsavebox{\mathbox}\newsavebox{\mathquote}
\makeatletter
\newcommand{\mq}[1]{% \mathquotes{<stuff>}
  \savebox{\mathquote}{\text{"}}% Save quotes
  \savebox{\mathbox}{$\displaystyle #1$}% Save <stuff>
  \raisebox{\dimexpr\ht\mathbox-\ht\mathquote\relax}{"}#1\raisebox{\dimexpr\ht\mathbox-\ht\mathquote\relax}{''}
}
\makeatother

% Paragraph mit Zähler (Section-Weise)
\newcounter{cparagraphC}
\newcommand{\cparagraph}[1]{
\stepcounter{cparagraphC}
\paragraph{\thesection{}-\thecparagraphC{} #1}
%\addcontentsline{toc}{subsubsection}{\thesection{}-\thecparagraphC{} #1}
\label{\thesection-\thecparagraphC}
}
\makeatletter
\@addtoreset{cparagraphC}{section}
\makeatother


% (Vorlesungs-)Folien einbinden:
\makeatletter
% Folien-Pfad (inkl. Dateiname ohne Endung und ggf. ohne Nummerierung)
\newcommand{\@slidePath}{}
\newcommand{\slidePath}[1]{\renewcommand{\@slidePath}{#1}}
% Folien Eigenschaften
\newcommand{\@slideScale}{0.5}
\newcommand{\slideScale}[1]{\renewcommand{\@slideScale}{#1}}
\AtBeginDocument{
% Folien von einer Datei skaliert
\newcommand{\slide}[2][\@slideScale]{\slides[#1]{}{#2}}
\newcommand{\slideTrim}[6][\@slideScale]{\slides[#1 , clip,  trim = #5cm #4cm #6cm #3cm]{}{#2}}
% Folien von mehreren nummerierten Dateien skaliert
\newcommand{\slides}[3][\@slideScale]{\begin{center}
\includegraphics[page=#3, scale=#1]{\@slidePath #2.pdf}
\end{center}}
}
\makeatother

% \emph{} anders definieren
\makeatletter
\DeclareRobustCommand{\em}{%
  \@nomath\em \if b\expandafter\@car\f@series\@nil
  \normalfont \else \scshape \fi}
\makeatother

% unwichtiges
\newcommand{\unimptnt}[1]{{\transparent{0.5}#1}}

% alph. enumerate
\newenvironment{anumerate}{\begin{enumerate}[label=(\alph*)]}{\end{enumerate}} % Alphabetische Aufzählung

% Hanging parameters
\newcommand{\hangpara}[1]{\par\noindent\hangindent+2em\hangafter=1 #1\par\noindent}

%% EINFACHE BEFEHLE

% Abkürzungen Mathe
\newcommand{\EE}{\mathbb{E}}
\newcommand{\QQ}{\mathbb{Q}}
\newcommand{\RR}{\mathbb{R}}
\newcommand{\CC}{\mathbb{C}}
\newcommand{\NN}{\mathbb{N}}
\newcommand{\ZZ}{\mathbb{Z}}
\newcommand{\PP}{\mathbb{P}}
\renewcommand{\SS}{\mathbb{S}}
\newcommand{\cA}{\mathcal{A}}
\newcommand{\cB}{\mathcal{B}}
\newcommand{\cC}{\mathcal{C}}
\newcommand{\cD}{\mathcal{D}}
\newcommand{\cE}{\mathcal{E}}
\newcommand{\cF}{\mathcal{F}}
\newcommand{\cG}{\mathcal{G}}
\newcommand{\cH}{\mathcal{H}}
\newcommand{\cI}{\mathcal{I}}
\newcommand{\cJ}{\mathcal{J}}
\newcommand{\cM}{\mathcal{M}}
\newcommand{\cN}{\mathcal{N}}
\newcommand{\cP}{\mathcal{P}}
\newcommand{\cR}{\mathcal{R}}
\newcommand{\cS}{\mathcal{S}}
\newcommand{\cZ}{\mathcal{Z}}
\newcommand{\cL}{\mathcal{L}}
\newcommand{\cT}{\mathcal{T}}
\newcommand{\cU}{\mathcal{U}}
\newcommand{\cX}{\mathcal{X}}
\newcommand{\cV}{\mathcal{V}}
\renewcommand{\phi}{\varphi}
\renewcommand{\epsilon}{\varepsilon}
\renewcommand{\theta}{\vartheta}

% Verschiedene als Mathe-Operatoren
\DeclareMathOperator{\arccot}{arccot}
\DeclareMathOperator{\arccosh}{arccosh}
\DeclareMathOperator{\arcsinh}{arcsinh}
\DeclareMathOperator{\arctanh}{arctanh}
\DeclareMathOperator{\arccoth}{arccoth} 
\DeclareMathOperator{\var}{Var} % Varianz 
\DeclareMathOperator{\cov}{Cov} % Co-Varianz 

% Farbdefinitionen
\definecolor{red}{RGB}{180,0,0}
\definecolor{green}{RGB}{75,160,0}
\definecolor{blue}{RGB}{0,75,200}
\definecolor{orange}{RGB}{255,128,0}
\definecolor{yellow}{RGB}{255,245,0}
\definecolor{purple}{RGB}{75,0,160}
\definecolor{cyan}{RGB}{0,160,160}
\definecolor{brown}{RGB}{120,60,10}

\definecolor{itteny}{RGB}{244,229,0}
\definecolor{ittenyo}{RGB}{253,198,11}
\definecolor{itteno}{RGB}{241,142,28}
\definecolor{ittenor}{RGB}{234,98,31}
\definecolor{ittenr}{RGB}{227,35,34}
\definecolor{ittenrp}{RGB}{196,3,125}
\definecolor{ittenp}{RGB}{109,57,139}
\definecolor{ittenpb}{RGB}{68,78,153}
\definecolor{ittenb}{RGB}{42,113,176}
\definecolor{ittenbg}{RGB}{6,150,187}
\definecolor{itteng}{RGB}{0,142,91}
\definecolor{ittengy}{RGB}{140,187,38}

% Textfarbe ändern
\newcommand{\tred}[1]{\textcolor{red}{#1}}
\newcommand{\tgreen}[1]{\textcolor{green}{#1}}
\newcommand{\tblue}[1]{\textcolor{blue}{#1}}
\newcommand{\torange}[1]{\textcolor{orange}{#1}}
\newcommand{\tyellow}[1]{\textcolor{yellow}{#1}}
\newcommand{\tpurple}[1]{\textcolor{purple}{#1}}
\newcommand{\tcyan}[1]{\textcolor{cyan}{#1}}
\newcommand{\tbrown}[1]{\textcolor{brown}{#1}}

% Umstellen der Tabellen Definition
\newcommand{\mpb}[1][.3]{\begin{minipage}{#1\textwidth}\vspace*{3pt}}
\newcommand{\mpe}{\vspace*{3pt}\end{minipage}}

\newcommand{\resultul}[1]{\underline{\underline{#1}}}
\newcommand{\parskp}{$ $\\}	% new line after paragraph
\newcommand{\parnl}{\unskip\parfillskip 0pt \par} % justified new line (brace line!)

% mathematische Abkürzungen
\newcommand{\corr}{\;\widehat{=}\;}
\newcommand{\mdeg}{^{\circ}}
\newcommand{\nok}[2]{\binom{#1}{#2}}	% n über k BESSER: \binom{n}{k}
\newcommand{\mtr}[1]{\begin{pmatrix}#1\end{pmatrix}}	% Matrix
\newcommand{\dtr}[1]{\begin{vmatrix}#1\end{vmatrix}}	% Determinante (Betragsmatrix)
\LetLtxMacro{\originalVec}{\vec}
\renewcommand{\vec}[1]{\underline{#1}}	% Vektorschreibweise
\newcommand{\imptnt}[1]{\colorbox{red!30}{#1}}	% Wichtiges
\newcommand{\intd}[1]{\,\mathrm{d}#1}
\newcommand{\diffd}[1]{\mathrm{d}#1}
% für Module-Rechnung: \pmod{}
\newcommand{\unit}[1]{\,\mathrm{#1}}
%\LetLtxMacro{\ntilde}{\tilde}
%\renewcommand{\tilde}{\widetilde}

% weiter Abkürzungen
\newcommand{\gdw}{genau dann wenn}
\newcommand{\lecdate}[1]{\marginpar{\small\textcolor{gray}{Vorlesung\newline #1}}}

\faculty{Fakultät Informatik/Mathematik}
\chair{Lehrstuhl für Programmierung, Compilerbau/Interpreter, Mikrocontroller}
\subject{lecture}
\title{Programmierung 2}
\author{Falk-Jonatan Strube}
\professor{Prof. Dr.-Ing. Arnold Beck}

\begin{document}
\maketitle
\tableofcontents

\section*{Hinweise}
Unterlagen unter:
\begin{lstlisting}[language=bash]
cd /home/rex/fi1/nestler/Programmierung_II_2016/
\end{lstlisting}
\paragraph{Compiler} 
\begin{itemize}
\item Intel i16, i13 (für Linux oder Visual Studio) \url{www.hocomputer.de} (kostenpflichtig)
\item gcc 5.3, 4.85 \url{gcc.gnu.org}
\end{itemize}
Zugriff auf Windows-Programme (Visual Studio 2013) in Linux-Laboren:
\begin{lstlisting}[language=bash]
rdesktop -f its56 # oder its59
\end{lstlisting}
Empfohlene Literatur: Breymann\cite{breymann2009c++}

\chapter{C++}
\section{Ein- und Ausgabe}
(siehe Folie CPP\_01\_stdio)
\begin{lstlisting}[language=C++]
#include <iostream>	// alternativ zu <stdio.h> in C
using namespace std;	// namespace: für bestimmte Abkürzungen (bpsw. cout anstatt std::cout)
// Hinweis: "::" zeigt, dass das davorstehende "static" ist (hier: std)

class integer {}; // class: Vergleichbar mit typedef

int main() {
	integer i0;	// i0: Instanz bzw. Objekt der Klasse integer
	cin.get();	// Eingabe abwarten
}
\end{lstlisting}
(vgl. integer.cpp)
\begin{lstlisting}[language=C++]
#include <iostream>
using namespace std;

class integer { // int - Variable in class verpacken
//	private:    // private ist default
		int i;  // this->i bzw. (*this).i
		// private nur für andere Klassen, andere Instanzen der gleichen Klasse können drauf zugreifen
	public: // wenn nichts steht, wird automatisch das vorherige angenommen. Hier: private
		integer(int i=0):i(i){  // Konstruktor und Defaultkonstruktor
			// i=0 default Wert, wenn keiner Angegeben
			// :i(i) übergebenes i wird dem i der Klasse zugewiesen: i_1(i_2) i_1 ist this->i, und i_2 ist das übergebene i
			cout<<"integer-Objekt i = "<<this->i<<endl; 
		}

		int get(){ return i; }

		void set(int i=0){ this->i = i; }

		// statische Methode: aufrufbar ohne Instanziierung der Klasse
		static integer add(integer i1, integer i2){  // Wertkopien von i1 und i2
			// return integer(i1.i + i2.i); // alternativ und explizit: Konstruktor-Aufruf
			// return erstellt eine Kopie (mit Konstruktur erstellt)!
			// return i1.get()+i2.get();    // Umwandlung int nach integer, Aufruf Konstruktor implizit
			return i1.i + i2.i;
			// i1.i Möglich, da innerhalb der Klasse integer und somit privates i sichtbar
		}
};

auto max(int x, int y) -> int { return x>y ? x : y; }	// Lambda-Funktion
// auto: Rückgabetyp ergibt sich aus dem Kontext bzw. über das "-> int"

template<typename Typ1, typename Typ2>	// Weiterentwicklung Makro: wählt automatisch Typ aus
auto quotient(Typ1 a, Typ2 b) -> decltype(a/b) { return a/b; }

auto main() -> int {
	auto k = 0;        // C++11: da 0 vom Typ int ist auch k vom Typ int
	decltype(k) j = 5; // C++11: da k vom Typ int ist auch j vom Typ int
	char *c = nullptr; //C++ 11: Zeigerliteral
	int *ip = NULL;

	integer i0(5), i1=4;     // 2 (alternative) Initialisiierungen von Objekten
	// i1=4 nur möglich, wenn 1 Parameter gefordert ist.
	cout<<"i0.i = "<<i0.get()<<endl;
	// cout im Vergleich zu printf() typsicher.
	cout<<"i0.i + i0.i = "<<integer::add(i0, i0).get()<<endl; // Aufruf static-Methode add
	integer i3 = integer::add(i0, i0);                        // Initialisierung von i3
	cout<<"i3.i     = "<<i3.get()<<endl;
	i0.set(22);
	cout<<"i0.i     = "<<i0.get()<<endl;
	cout<<"max(3,5) = "<<max(3,5)<<endl;
	cout<<"5/3      = "<<quotient(5, 3)<<endl;
	cout<<"5.0/3.0  = "<<quotient(5.0, 3.0)<<endl;
	cout<<"b / 1    = "<<quotient('b', '1')<<endl;
	cin.get();
}
\end{lstlisting}

(vgl iostream.pdf)
\begin{itemize}
\item nach jeder cin Eingabe: „cin.clear();“, damit Fehler ignoriert werden um weiter cin's abhandeln zu können (vgl. robust\_ea)
\end{itemize}
Einlesen:
\begin{lstlisting}[language=C++]
char sc;
cout << "sc=";
cin >> sc;
cin.clear();	// clear, um die Eingabezeile freizumachen, damit man nicht an Falscher Eingabe hängen bleibt
cin.ignore(INT_MAX, '\n');	// braucht #include <limits.h>
cout << "sc" << dec << (int) sc << endl;

// alternativ:
char vb[128];
cout << "s=";
cin.getline(vb, sizeof(vb), '\n');	// lesen als String, dann wieder umwandeln (liest auch Leerzeichen ein)
sc = atoi(vb);	// braucht #include <cstdlib>

// alternative zu getline:
cin.get(...);	// lässt aber \n im Strom
cin.get();

// alternativ
cin >> setw(sizeof(vb)) >> vb;	// verhindert Überlauf
sc = atoi(vb);

// alternativ
String s;	// braucht #include <string>
size_t ie=0;
cin >> s;
unsigned int ni = stoi(s, &ie, 10);

// alternativ
getline(cin, s, '\n');
double d = stod(s, &ie);

// zum compilieren: g++ p2a1.cpp -std=c++11 -o a.out
\end{lstlisting}

robust\_ea1.cpp:
\begin{lstlisting}[language=C++]
	do {
		cout<<"d = "; cin>>d; 	// einlesen
		if(cin.eof()) break;		// break bei Strg+D oä.
		if(cin.fail() || (cin.peek() != '\n')){		// ist nächstes Zeichen ungültig?
			cin.clear(); cin.ignore(INT_MAX,'\n');	// Strom zurücksetzen und zum \n gehen
			continue;
		}
		break; % Schleife verlassen, wenn korrekte Eingabe
	} while(true);
	
	if(cin.eof()){ cin.clear(); cout<<"eof\n"; }
	else {
		cin.clear(); cin.ignore(INT_MAX,'\n');
		cout<<"Wert d = "<<d<<endl;
	}
	cin.ignore();
\end{lstlisting}

\section{Defaultargumente}
Defaultargumente müssen immer von rechts angefangen definiert sein:
\begin{lstlisting}[language=C++]
myFunc(int i = 5, int j = 7)	// korrekt
myFunc(int i, int j = 7)			// korrekt
myFunc(int i = 5, int j)			// falsch!!!
\end{lstlisting}

\section{Überladen}
overload.pdf

Hinweis: cast auf zwei Möglichkeiten:
\begin{lstlisting}[language=C++]
int i = 5;
double d;
d = (long) i;
d = long(i);
\end{lstlisting}

\section{Typisierte Konstanten}

\section{Referenzen}
referenzen.pdf\\
Ein Speicherplatz wird mit mehreren Variablen-Namen beschrieben.

\section{String0}
\section{String}
\section{const}
const\_mutable.pdf\\
Faustregel: Alle Funktionen, die nichts ändern, immer das const anfügen.
\section{Folge}
Initfolge/Initfolge.pdf
\section{Kopierkonstrukturen}
Kopierkonstruktoren.pdf

\section{Friend}
Zugriff auf private Klassen von außerhalb:
\begin{itemize}
\item Funktionen
\item Alle Methoden anderer Klassen
\item spezifische Methoden anderer Klassen
\end{itemize}
Muss von Klasse mit privater Funktion festgelegt werden. Freundschaft wird nicht „erwidert“, ist nicht reflexiv (bsp. friendreflex).

\section{Vererbung}
vererben.pdf

\section{Virtuelle Funktionen}
\begin{itemize}
\item überschriebene Funktionen 
\item virtuelle Funktionen $\Rightarrow$Späte Bindung
\item rein virtuelle Funktionen\\
eine Klasse mit rein virtuellen Funktionen ist abstrakt
\item VMT virtual method table ABB 124
\end{itemize}
besch1.cpp (Beschäftigte)
\begin{lstlisting}[language=C++]
#include <iostream>
#include <string>
using namespace std;

class Besch{
	private:
		string name;
	public:	
		Besch(string name): name(name){}	// :name(name) --> Member-Initialisierer
		string getName() {return name;}
		void setName(string name) {this->name = name;}
		virtual void display(ostream& os) {os<<name;}	// virtuelle Funktion
		virtual double calc(){return 0.0;}	// normalerweise Geld nicht mit double, sondern ganzzahlig rechnen!
		// wenn Funktion eigentlich nicht benötigt wird: Polymorphes Interface:
		// virtual double calc() = 0; // rein virtuelle Funktion
		// diese rein virtuelle Funktion macht den Datentyp Besch abstract! D.h. davon kann keine Instanz erzeugt werden (wie in der Main passiert)!
};
ostream& operator<<(ostream& os, Besch& b) {b.display(os);}

class Arbeiter:public Besch{	// public: durchsichtige Vererbung
	private:
		int stunden;
		double stdLohn;
	public:
		Arbeiter(string name, double stdLohn):Besch(name){ // :Besch(name) --> Basis-Initialisierer
			this->stdLohn = stdLohn;
		}
		int getStunden() { return stunden; }
		void setStunden(int stunden) { this->stunden = stunden; }	// usw. getter und setter
		void display(ostream& os){ Besch::display(); os<<""<<stdLohn<<" "<<stunden<<" "<<calc();}
		double calc(){ return stdLohn * stunden; }
};
ostream& operator<<(ostream& os, Arbeiter& a) { a.display(os); }

class Angest:public Besch{
	private:
		double gehalt;
	public:
		Angest(string name, double stdLohn):Besch(name){
			this->gehalt = gehalt;
		}
		void display(ostream& os){ Besch::display()<<" "; os<<gehalt<<" ";}
		double calc(){ return gehalt; }
};
ostream& operator<<(ostream& os, Angest& a) { a.display(os); }

class Haendler:public Arbeiter{
	private:
		double prov;
		double ums;
	public:
		Haendler(string name, double stdLohn, double prov):Arbeiter(name, stdLohn){this->prov=prov;};
		void setUms(double ums){this->ums = ums;}
		double calc(){return Arbeiter::calc() ++ ums*prov;}
		void display(ostream &os){Arbeiter::display(os);}
}

int main() {
	Besch* b1 = new Besch("Hans Huckebein");
	b1->display(cout); cout<< " Euro"<<endl;
	cout<< (*b1) << " Euro" << endl;
	cout<< "=============="<<endl;
	Arbeiter* b2 = new Arbeiter("Moriz Lehmann",  8.99);
	b2->setStunden(140);
	b2->display(cout); cout<< " Euro"<<endl;
	cout<< (*b2) << " Euro"<< endl;
	cout<< "=============="<<endl;
	Angest* b3 = new Arbeiter("Friedrich Lempel",  3099.00);
	b3->display(cout); cout<< "Euro"<<endl;
	cout<< (*b3) << " Euro"<< endl;
	cout<< "=============="<<endl
	Haendler* b4 = new Handler("Bang Ohlufson",  15.80, 0.1);
	b4->setStunden(350);
	b4->setUms(12000);
	b4->display(cout); cout<< "Euro"<<endl;
	cout<< (*b3) << " Euro"<< endl;
	cout<< "=============="<<endl
	
	Besch* be[] = {b1,b2,b3,b4};
	for (int i=0; i<4; i++){
		cout<<i<<": "; 
		be[i]->display(cout); cout<<endl;	// gibt 0 für alle zurück, da nur calc() (und display()) vom Typ Besch ausgeführt, da das Array aus Besch besteht! Die Funktion calc() wurde überschrieben (Vgl. Überladen, wenn bei gleichen Namen verschiedene Paramater angegeben sind)!
		// Pointer im Array bestimmt die ausgeführten Funktionen!
		// Lösung: "virtual" als Schlüsselwort vor Funktion (s.o.). Angeben bei erstem Auftauchen in Vererbungshierarchie
		cout<<(*(be[i]))<<endl<<"+++++++++++++++"<<endl;	// geht auch, da im Operator (die ja nicht virtuell gemacht werden kann, weil sie kein Member ist) die virtuelle display-Fkt aufgerufen wird.
	}	
	
	return 0;
}
\end{lstlisting}
Vererbungs Hierarchie:\\
ABB 125\\
Potentielle Mehrfachvererbung (bei Manager):
\begin{itemize}
\item Hat mehrfache Namen (wegen unterschiedlichen Vererbungslinien)
\end{itemize}
Problemlösung: Arbeiter, Angest und Händler von virtual Besch erben lassen (dort, wo es sich aufzweigt):
\begin{lstlisting}[language=C++]
class Arbeiter: virtual public Besch{};
\end{lstlisting}
Hat immer noch Einschränkungen!

\section{Templates}
Templates Folie (enthält auch exception).
Vor allem für Containerklassen/Collections verwendet (Vektoren, Hashtables, Listen).\\
Schafft Möglichkeit Klassen parameterbehaftet zu generieren (generische Klassen in Java).\\
Funktionen von Templates müssen (im Gegensatz zur bisherigen Vorgehensweise) im header-File ausprogrammiert werden. Alles, was zum Template gehört kommt ins header-file!!!
\paragraph{Beispiel Array} Intarr.h
\begin{lstlisting}[language=C++]
const int SizeArr = 24;

// Konstruktor, Kopierkonstruktor und Destruktor werden bei Klassen mit Pointer immer benötigt!
class IntArr{
	public:
		IntArr (int Gr = SizeArr ) 	// Konstruktor
		IntArr (const IntArr&) 			// Kopierkonstruktor
		~IntArr() (delete IA;}			// Destruktor
		IntArr & operator = (const IntArr&);
		int& operator [] (int i);	
		int getNum() const {return Size;}
		void resize(int NewSz);
	private:
		int Size;
		int * IA;
};
\end{lstlisting}
Auszug Intarr.cpp
\begin{lstlisting}[language=C++]
...
IntArr :: IntArr (int Sz){
	Size = Sz;
	IA = new int [Size];
	for (int i=0; i<Sz; ++i) IA[i] = 0;
}
\end{lstlisting}
Modifikation um es generisch zu machen (Template), Inhalte von Intarr.cpp müssen eingefügt werden:\\
Intarr.h
\begin{lstlisting}[language=C++]
const int SizeArr = 24;

template <class T>
class IntArr{
	public:
		IntArr (int Gr = SizeArr );
		IntArr (const IntArr<T>&);
		~IntArr() (delete IA;}
		IntArr & operator = (const IntArr<T>&);
		T& operator [] (int i);	
		int getNum() const {return Size;}
		void resize(int NewSz);
	private:
		int Size;
		T * IA;
};

//... + Inhalt von Intarr.cpp!!!
template <class T>
IntArr<T> :: IntArr (int Sz){
	Size = Sz;
	IA = new T [Size];
	for (int i=0; i<Sz; ++i) IA[i] = 0;
}

template <class T>
IntArr<T> :: IntArr (const IntArr<T>& Other){
	Size = Other.Size;
	IA = new T[Size];
	for (int i=0; i<Size; i++) IA[i]=Other.IA[i];
}
usw.
\end{lstlisting}
in der Main dann:
\begin{lstlisting}[language=C++]
// vorher: IntArr IA (10);
IntArr<int> IA (10);
// oder auch 
IntArr<double> IA (10);
// auch möglich 
IntArr<Fraction> IA (10); // wobei Fraction in einer anderen Klasse definiert ist.
\end{lstlisting}

\paragraph{Beispiel Listen}
cobject.h

\chapter{Java}
\section{Grundlagen}
Java0intro.pdf\\
Klassenname und Dateiname müssen exakt übereinstimmen. Jede Datei darf nur eine Klasse enthalten.
\begin{lstlisting}[language=Java]
public class HelloOne{	// Dateiname: HelloOne.java
	public static void main(String args[]}{	// im Vergleich zu C++ gibt es keinen public-Bereich, die Eigenschaft muss vor jede Funktion/Variabel geschrieben werden.
		System.out.println("welcome to java");
		System.out.println("today is " + new java.util.Date());
		System.out.printf("JAVA test: %d\n", 26);
	}
}
\end{lstlisting}
javac HelloOne.java zum Übersetzen\\
java HelloOne zum Ausführen
\begin{lstlisting}[language=Java]
public class HelloEcho{
	public static void main(String args{]){
		for (int i = 0; i<args.length; i++)
			System.out.println(args[i]);
		// oder:
		for (String s:args)
			System.out.println(s);
		System.exit(0);
	}
}
\end{lstlisting}
\subsection{Parsen von Daten}
\begin{lstlisting}[language=Java]
public class HelloParse{
	public static void main(String args{]){
		int i;
		double d;
		i = Integer.parseInt(args[0]);
		d = Double.parseDouble(args[1});
		System.out.println("i: "+i+" d: "+d);
	}
}
\end{lstlisting}
\subsection{GUI mit AWT}
\begin{lstlisting}[language=Java]
import java.awt.*;
import java.awt.event.*;
import java.util.*;
public class HelloTwo extends Panel {	// extends: erweitert. Also: "erbt von"
	Label l1=new label("welcome to java", Label.CENTER);
	Label l2=new Label("today is "+new Date(), Label.CENTER);	// wegen import javo.util.* kann hier nur Date() geschrieben werden.
	Button close = new Button ("Close");
	
	public HellowTwo(){
		setFont(new Font("System", Font.PLAIN, 24));
		setLayout(new GridLayout(3,1,0,10));
		add(l1);	// alternativ: this.add(l1);
		add(l2);
		close.addActionListener(new ActionListener(){
			public void actionPerformed(ActionEvent e){
				// bpsw erweiterbar mit System.out.printl("Close...");
				System.exit(0);
			}
		});
	}
	
	public static void main(String args[]){
		Frame f=new Frame("first Desktop Application");
		HelloTwo h = new HelloTwo();
		f.add(h);
		f.setVisible(true);
		f.pack();
	}
}
\end{lstlisting}

\subsection{Applet}
\begin{lstlisting}[language=Java]
import java.applet.*;
import java.awt.*;

public class HelloThree extends Applet{
	public void init(){
		Panel h= new HelloTwo();
		add(h);
	}
}
\end{lstlisting}
\begin{lstlisting}[language=HTML]
<html>
<body>
<applet code=HelloThree width=940 height=540></applet>
</body>
</html>
\end{lstlisting}
Zu jedem Applet gehört ein dazugehöriges HTML-File. Damit kann das Applet über einen Applet-Viewer benutzt werden. Hier funktioniert der Close-Button aber nicht. 
\paragraph{Angepasstes HelloTwo} Hiermit wird der Close-Button implementiert, sodass er im Applet nicht erscheint.
\begin{lstlisting}[language=Java]
// Quit-Button wird erst in main-Funktion eingeführt
\end{lstlisting}
java1.pdf
\subsection{Operatoren}
\begin{itemize}
\item $>>$ schiebt nach rechts und fügt abhängig vom ersten Bit (Vorzeichen) $0$en (1. Bit $0$) oder $1$en (1. Bit $1$) nach. $>>>$ schiebt immer $0$en nach.
\end{itemize}

\subsection[Eingabekonvertierung, Formatierte Ausgabe]{Eingabekonvertierung, Formatierte Ausgabe / Konvertierung numerischer Werte}
\begin{lstlisting}[language=Java]
public class Num{
	public static vaid main(String args[]){
		int x;
		x = Integer.parseInt(args[0]);	-- Exceptions bei keiner Eingabe und bspw. String-Eingabe
		System.out.printf("x: %d", x);
		System.out.println("x: "+x);
	}
}
\end{lstlisting}

\subsection{Vergleich von Objekten}
\begin{itemize}
\item == prüft Objekt gleich ist (nicht Inhalt!)
\item .equals() prüft Gleichheit von Objektinhalt
\end{itemize}
\begin{lstlisting}[language=Java]
public class Cmp{
	public static vaid main(String args[]){
		String s="Hans";
		if (args[0] == s)
			System.out.println("args[0]==s liefert true");
		else
			System.out.println("args[0]==s liefert false");
		if (s.equals(args[0]))
			System.out.println("s.equals(args[0]) liefert true");
		else
			System.out.println("s.equals(args[0]) liefert false");
		// ergibt das erste Mal false (weil Objekte ungleich) und das zweite Mal true (weil Objektinhalt gleich).
			
		if ("Hans" == s)
			System.out.println("Hans==s liefert true");
		else
			System.out.println("Hans==s liefert false");
		if (s.equals("Hans"))
			System.out.println("s.equals(\"Hans\") liefert true");
		else
			System.out.println("s.equals(\"Hans\") liefert false");
		// ergibt beides mal true, weil der Compiler optimiert und der "Hans" String nicht doppelt erzeugt wurde.
	}
}
\end{lstlisting}

\subsection{Arrays}
\begin{lstlisting}[language=Java]
int intArray[];	// [] bleibt hier immer leer!
intArray = new int[2]	// jetzt [] füllen
// alternativ, inkl. Initialisierung:
int intArray[] = {1,2,3};

// mehrdimensional:
int mda[][] = { {1,2}, {1,2,3}, {1,2,3,4,5} };
// alternativ:
int mda[][];
mda = new int{2][];
mda[0] = new int[2];
mda[1] = new in[5];
// oder auch:
mda = new int[2]{10];
\end{lstlisting}
Vgl. Byte Array (zum Verschicken über Netzwerk: Serialisierung)

\subsection{Statements}
if- und Schleifenbedingung braucht unbedingt boolsche Werte.\\
Also bspw. nicht "if(i)", sondern "if(i!=0)".

\begin{lstlisting}[language=Java]
public class Cmp{
	public static void main(String args[]){		
		for ( String x:args) 	// for-each alternativ zu for(int i=0; i<args.length; i++) {x=args[i]; System.out.println(x);}
			System.out.println(x);
	}
}
\end{lstlisting}

\subsection{Exceptions}
\begin{lstlisting}[language=Java]
class tryDemo{
	...
}
\end{lstlisting}

\section{Klassen}
Java2\_classes.pdf

\begin{lstlisting}[language=Java]
class LocalVar{
	public static void main(String args[]){
		int i,j;
		j=i+1;	// ergibt Fehler, weil i noch nicht initialisiert wurde.
		System.out.println(j);
	}
}
\end{lstlisting}

\subsection{Initialisierer}
Funktion ohne Funktionskopf. Wird ausgeführt, wenn Instanz erstellt wird. Ein static-Initialisierer wird ausgeführt, sobald die Klasse geladen wird.
\subsection{Vererbung}
Überschriebene Methoden in Java verhalten sich wie virtuelle Methoden in C++.

\begin{lstlisting}[language=Java]
import java.awt.*;
import java.awt.event.*;

public class myPanel extends Component implements WindowListener{
	private String myString;
	
	public void windowActivated(WindowEvent e){}
	public void windowClosed(WindowEvent e){}
	public void windowClosing(WindowEvent e){System.exit(1);}
	public void windowDeactivated(WindowEvent e){}
	public void windowDeiconified(WindowEvent e){}
	public void windowIconified(WindowEvent e){}
	public void windowOpened(WindowEvent e){}
	
	@Override
	public Dimension getPreferredSize()	{
		return new Dimension(400,200);
	}
	
	myPanel(String s){
		myString =s;
	}
	
	public static void main(String args[]){
		Frame f=new Frame (args.length>0?args[0]:"no String - no fun");
		myPanel p=new myPanel(args.length>0?args[0]:"no String - no fun");
		System.out.println(p.myString+" in main");
		f.add(p);
		f.addWindowListener(p);
		f.pack(); //setSize(400,200);
		f.setVisible(true);
		f.repaint();
	}
	
	@Override
	public void paint(Graphics gc){
		//System.out.println(myString+" in paint");
		gc.drawString(myString,10,20);
	}
}
\end{lstlisting}

\section{Abstract Windowing Toolkit (AWT)}
java4\_awt.pdf
% http://www.informatik.htw-dresden.de/~beck/java/Vorlesung/Beispielprogramme/AWT/Muster.java AUF SPICKER!
\subsection{BorderLayout}
Sollen mehrere Buttons in einen Bereich (bspw. NORTH), muss ein extra Panel dort eingefügt werden:
\begin{lstlisting}[language=Java]
...
Panel pTop = new Panel(new FlowLayout());
poben.add(b1);
poben.add(b2);
add(poben,BorderLayout.NORTH);
...
\end{lstlisting}

\subsection{CardLayout}
Mehrere Karten, man sieht immer nur eine. Wechseln mit next, previous oder show.

\paragraph{Mehrere Buttons usw. hinzufügen} $ $
\begin{lstlisting}[language=Java]
String[] myLabels = {"Max", "Moritz", "Bolte"};
...
Panel p=new Pannel(new FlowLayout());
for (int i=0; i<myLabels.length; i++){
	Button b = new Button (myLabels[i]);
	p.add(b);
}
this.add(p.BorderLayout.CENTER);
\end{lstlisting}

\section{Eventhandling}
java4\_eventhandlig.pdf

\subsection{Inner classes}
4 Arten merken!

\section{Java I/O}
java3\_io.pdf

\subsection{HexDump}


\chapter{Prüfungsinhalt}
C++: manipulator, cout, <<, Ausgabestring formatieren, Operatorüberladung, Vererbung, Überladen/überschriebene/virtuelle Funktionen//
Java bis einschließlich AWT

\newpage
\printbibliography
\end{document}